\chapter{הבעיה וההזדמנות}\label{ch:problem}

\section{גגות שאינם בשימוש וערי חום}
בערים צפופות, משטחי בטון ואספלט סופגים חום ומאיצים את תופעת אי החום העירוני. סקירות מקצועיות מראות שגגות ירוקים וגגות מחזירי אור מקטינים את טמפרטורת המשטח ומשפרים את תנאי הנוחות בעיר \cite{santamouris2014}. מעבר לקירור, גגות ירוקים מספקים שירותי מערכת אקולוגית נוספים, כמו ניקוז מי גשמים ובידוד תרמי \cite{oberndorfer2007}.

ברוב בנייני המגורים בגוש דן הגג נשאר ריק או משמש לציוד טכני. אין לו הכנסה, ועל הוועד מוטלת עלות תחזוקה. גג שמגדלים עליו ירקות הופך מנטל לנכס: הוועד מקבל \RRrentShare\% מההכנסות, והבניין מקבל שטח ירוק מקרר.

\section{מרחק, טריות ושרשרת אספקה}
ירקות עלים הם מוצר מתכלה במיוחד. הם נקטפים בשדה, עוברים מיון, אריזה, הובלה ומחסן, ומגיעים למסעדה ימים אחרי הקטיף. כל שלב מקצר את חיי המדף ומגדיל פחת. מסעדה שמזמינה עלים טריים מגג במרחק של פחות מק״מ מקבלת מוצר שנקטף באותו בוקר ואינה משלמת על קירור והובלה ארוכה.

עבור המסעדן המשמעות ברורה: פחות פחת, טעם עקבי יותר, וסיפור מותג של מזון מקומי. עבור העיר: פחות משאיות ופחות פליטות. הערכות אלה הן הנחות עבודה של הצוות, שנבדקות בתוכנית הפיילוט, ואינן נתון סטטיסטי.

\section{מחסור במים}
ישראל היא מדינה עם משאבי מים מוגבלים, ובה חקלאות מסורתית צורכת כמויות גדולות של מים להשקיה. הידרופוניקה ממחזרת את תמיסת ההזנה במעגל סגור, וכך צריכת המים לק״ג תוצרת נמוכה בהרבה משדה פתוח. ההשוואה הכמותית מפורטת בפרק~\ref{ch:solution}, ובעיקר באיור~\ref{fig:resources}, המבוסס על \cite{barbosa2015}.

\section{למה עכשיו}
שלושה כוחות מתכנסים.

\begin{enumerate}
  \item \textbf{ביקוש עולמי למזון בר-קיימא.} דוח ה-\en{FAO} מצביע על הצורך להגדיל את ייצור המזון תוך שימוש יעיל יותר במשאבים \cite{fao2017}.
  \item \textbf{בשלות טכנולוגית של חקלאות סביבה מבוקרת.} חקלאות משולבת בניינים ושליטה מבוססת חיישנים נסקרות בספרות כאפשרות בת-קיימא \cite{specht2014}. לפי הערכת הצוות הן זולות ואמינות מספיק לשימוש מסחרי בקנה מידה קטן, וההערכה נבדקת בפיילוט.
  \item \textbf{סוכני בינה מלאכותית.} מודלי שפה יכולים לפעול בלולאה של הנמקה ופעולה \cite{yao2023}. הנחת העבודה שלנו היא שגישה כזו תאפשר לנהל מערך מבוזר של גגות עם מגדל אנושי אחד לכמה גגות, והפיילוט יבדוק אותה.
\end{enumerate}

\begin{insightbox}[למה זה הזמן]
הנכס (גגות) קיים, הביקוש הצפוי (מסעדות מקומיות) ייבדק בפיילוט, והחוליה החסרה הייתה ניהול זול ומבוזר. \en{RootOS} סוגר את הפער הזה, ולכן פרק~\ref{ch:solution} מתאר אותו בפירוט.
\end{insightbox}
